\documentclass[12pt]{article}

\setlength{\parindent}{0pt}
\setcounter{secnumdepth}{3}
\usepackage[margin=1in]{geometry}

% packages
\usepackage{amsmath, amsfonts, amsthm, amssymb}
\usepackage{graphicx, float}
\usepackage{mathtools}
\usepackage{titlesec}
\usepackage{interval}
\usepackage{hyperref}
\usepackage{siunitx}
\usepackage{titling}
\usepackage{vwcol}
\usepackage{setspace}
\usepackage{empheq}
\usepackage{cancel}
\usepackage{esdiff}
\usepackage{multicol}
\usepackage{mdframed}
\usepackage{esdiff}
\usepackage{tikzsymbols}
\usepackage{multicol}
\usepackage{tikz}
\usepackage{varwidth}
\usepackage{pgfplots}
\usepackage{fancyhdr}
\usepackage{enumitem}
\usepackage{tcolorbox}

\intervalconfig {
	soft open fences
}

% header
\pagestyle{fancy}
\fancyhf{}
\lhead{Ethan Fan}
\rhead{AP Calculus BC}
\cfoot{\thepage}

% section format
\titleformat{\section}
  {\large \bfseries}
  {}
  {0em}
  {}
\titlespacing*{\section}{0pt}{0.5em}{0.3em}

\titleformat{\subsection}[runin]
  {\bfseries}
  {}
  {0em}
  {}
\titlespacing*{\subsection}{0pt}{0pt}{0em}

\titleformat{\subsubsection}[runin]
  {\itshape}
  {(\roman{subsubsection})}
  {0em}
  {}

% commands
\newcommand{\alignedintertext}[1]{%
  \noalign{%
    \vskip\belowdisplayshortskip
    \vtop{\hsize=\linewidth#1\par
    \expandafter}%
    \expandafter\prevdepth\the\prevdepth
  }%
}

\newcommand*{\paren}[1]{\ensuremath\left(#1\right)}
\newcommand*{\limit}[2][x]{\ensuremath{\displaystyle\lim_{#1\to#2}}}
\newcommand*{\Limit}[3][x]{\ensuremath{\displaystyle\lim_{#1\to#2}\left[#3\right]}}
\newcommand*{\deriv}[1][x]{\ensuremath{\dfrac{\mathrm{d}}{\mathrm{d}#1}}}
\newcommand*{\Deriv}[2][x]{\ensuremath{\dfrac{\mathrm{d}}{\mathrm{d}#1}\left[#2\right]}}
\newcommand{\ddx}{\dfrac{d}{dx}}
\newcommand{\dydx}{\dfrac{dy}{dx}}
\newcommand{\dydt}{\dfrac{dy}{dt}}
\newcommand{\dxdt}{\dfrac{dx}{dt}}
\newcommand*{\iinteg}[2][x]{\ensuremath{\displaystyle\int #2\;\mathrm{d}#1}}
\newcommand*{\dinteg}[4][x]{\ensuremath{\displaystyle\int_{#2}^{#3}#4\;\mathrm{d}#1}}
\newcommand*{\abs}[1]{\ensuremath{\left|#1\right|}}
\newcommand*{\eps}{\varepsilon}
\newcommand*{\real}{\ensuremath\mathbb{R}}
\newcommand*{\integer}{\ensuremath\mathbb{Z}}
\newcommand*{\posint}{\ensuremath\mathbb{Z^+}}
\newcommand*{\cbrt}[1]{\ensuremath{\sqrt[3]{#1}}}
\newcommand*{\lheqzero}{\ensuremath{\underset{\text{L'H}}{\overset{\left[\frac00\right]}{=}}}}
\newcommand*{\lheqinfty}{\ensuremath{\underset{\text{L'H}}{\overset{\left[\frac{\infty}{\infty}\right]}{=}}}}
\newcommand{\tor}{\text{ or }}
\newcommand{\floor}[1]{\lfloor #1 \rfloor}

\newcommand{\vem}{\vspace{0.5em}}
\newcommand{\example}[1]{\section{Example #1}}
\newcommand{\problem}[1]{\section{Problem #1}}
\newcommand{\subproblem}[1]{\subsection{(#1)} \,}

\DeclareMathOperator{\dne}{DNE}
\DeclareMathOperator{\sgn}{sgn}

\newcommand{\definition}[1]{\begin{tcolorbox}[colback=red!5!white,colframe=red!75!black,parbox=false] #1 \end{tcolorbox}}

\newcommand{\theorem}[2]{\begin{tcolorbox}[title={#1},colback=blue!5!white,colframe=blue!75!black,parbox=false] #2 \end{tcolorbox}}

\allowdisplaybreaks
% \postdisplaypenalty=100000

% document
\begin{document}

\vspace*{0cm}

% opening
\begin{center}
    {\Large \bfseries Problem Set 20} \\[0.5cm]
    {\large Volumes of Revolution I: Disks and Washers} \\[0.35cm]
    {\normalsize October 26, 2026}
\end{center}

% problems

\problem{1}

\subproblem{b}
\begin{gather*}
    V=\pi \int_{r-h}^{r} r^2-x^2\, dx = \pi (r^2 x-\frac{x^3}{3})\Big|_{r-h}^{r}\\
    =\pi (r^3-\frac{r^3}{3} - r^3 + r^2h + \frac{r^3}{3}-r^2h+rh^2-\frac{h^3}{3})\\
    = \boxed{\pi h^2(r-\frac{h}{3})}
\end{gather*}

\problem{2}

\subproblem{b}
\begin{gather*}
    A=\pi kx\\
    V=\pi \int_0^ a kx \, dx=\pi (\frac{1}{2}kx^2)\Big|^a_0= \boxed{\frac{\pi}{2}ka^2} 
\end{gather*}

\problem{4}

\subproblem{a}
\begin{gather*}
    V=\pi \int_1^2(2-\frac{x}{2})^2\, dx = \pi \int_1^2 \frac{x^2}{4}-2x+4\, dx\\
    =\pi (\frac{x^3}{12}- x^2+4x)\Big|_1^2 =\pi (\frac{7}{12}-3+4)=\boxed{\frac{19\pi}{12}}
\end{gather*}

\subproblem{b}
\begin{gather*}
    V=\pi \int_1^2 \frac{1}{x^2}\, dx=\pi (-\frac{1}{x})\Big|_1^2=\boxed{\frac{\pi}{2}}
\end{gather*}

\subproblem{c}
\begin{gather*}
    V=\pi \int_0^9 4y\, dy =2\pi (y^2)\Big|^9_0=\boxed{162\pi }
\end{gather*}

\subproblem{d}
\begin{gather*}
    V=\pi \int_0^1x^2-x^6\, dx =\pi (\frac{x^3}{3}-\frac{x^7}{7})\Big|^1_0= \boxed{\frac{4\pi}{21}}
\end{gather*}

\subproblem{f}
\begin{gather*}
    V=\pi \int_0^4 (2-\frac{y}{2})^4\, dy = \begin{bmatrix}
        u=2-\frac{y}{2}\\
        -2du=\, dy
    \end{bmatrix}
    \pi \int_0^2 2u^4\, dy= 2\pi (\frac{u^5}{5})\Big|^2_0=\boxed{\frac{64\pi}{5}}
\end{gather*}

\problem{6}
The true geometric area is expressed as $A=\pi (R^2-r^2)$.
\begin{gather*}
    V=\pi \int_0^1 (x^2-x^4)\, dx = \pi (\frac{x^3}{3}-\frac{x^5}{5})\Big|^1_0=\boxed{\frac{2\pi}{15}}\\
    R^2-2Rr+r^2=R^2-r^2\\
    2r(R-r)=0 \implies \boxed{r=0 \tor R}
\end{gather*}

\problem{7}

\subproblem{a}
\begin{gather*}
    V=\pi \int_0^{h } y\, dy=  \boxed{\frac{\pi h^2}{2}}
\end{gather*}

\subproblem{b}
\begin{gather*}
    V=\frac{\pi}{2}h^2\implies dV=\pi h \, dh\\
    \frac{dh}{dt}=\frac{dV}{dt}\cdot \frac{dh}{dV}=\frac{1}{4\pi}\cdot 2=\boxed{\frac{1}{2\pi} \text{ cm/s}}
\end{gather*}

\subproblem{c}
\begin{gather*}
    V(10)=50\pi  \\
    25\pi =\frac{\pi}{2}h^2\\
    h=\boxed{5\sqrt{2}\text{ cm}}
\end{gather*}
The bowl gets wider as the height increases.

\problem{8}

\subproblem{a}
\begin{gather*}
    (x-R)^2+y^2=a^2\\
    x=\boxed{R\pm \sqrt{a^2-y^2}}
\end{gather*}

\subproblem{b}
\begin{gather*}
    V=\pi \int_{-a}^a (R+ \sqrt{a^2-y^2})^2-(R- \sqrt{a^2-y^2})^2\, dy\\
    =\pi \int_{-a}^a 4R\sqrt{a^2-y^2} \, dy = \pi \cdot 4R\cdot \frac{1}{2}\pi a^2\\
    =\boxed{2\pi ^2Ra^2}
\end{gather*}

\subproblem{c}
\[
2\pi^2Ra^2=(2\pi R) (\pi a^2)
\]
$\pi a^2$: area of the rotated circle. \\
$2\pi R$: the distance traveled by the center of the circle. 

\problem{9}

\subproblem{a}
Upper bound: $y=\sqrt{r^2-x^2}$.\\
Lower bound: $y=\sqrt{r^2-\frac{h^2}{4}}$.

\subproblem{b}
\begin{gather*}
    V= \pi \int_{-\frac{h}{2}}^{\frac{h}{2}}  r^2-x^2-r^2+\frac{h^2}{4}\, dx\\
    = \pi ( - \frac{x^3}{3} +\frac{h^2}{4}x)\Big|_{-\frac{h}{2}}^{\frac{h}{2}} \\
    =2\pi\cdot \frac{2h^3}{24} = \boxed{\frac{\pi h^3}{6}}
\end{gather*}

\subproblem{c}
The volumes are the same. For the Earth, the radius of the hole is also extremely long, rendering the actual thickness extremely small, which cancels out and leaves an equal volume. \vem

\subproblem{d}
\begin{gather*}
    V=2\pi h^2(r-\frac{h}{3})\\
    =2\pi \cdot \frac{1}{4}\cdot \frac{5}{6}\\
    =\boxed{\frac{5\pi}{12}}
\end{gather*}
















\end{document}

